\section{Conclusion}\label{sec:conclusion}

Sixteen weeks of work produced three kinds of result. The first is
infrastructure: a single unit and harness in which rho, BSGS, the kangaroo
and every index-calculus variant are charged the same way on the same
one-target workloads, with sealed records, exact replays and admitted
intervals; and a ledger that refuses a claim with a missing field. The
second is a set of constant-factor improvements, the largest of which---the
closed-form roots of $\Sem3$---began as one sentence in a manuscript we were
asked to check, and ended as a proved identity, an oracle family covering
prime and binary fields, two to four summands, CPU and GPU, plain, tower
and dense-index bases, and the default decomposition strategy of the
codebase. The third is a map of what does not work, which is longer than
the list of what does: SAT that does not scale, symmetry that pays at
$n=17$ and costs at $n=31$, transfers that provably cannot change
decomposition counts, fibre-aware bases whose gap to rho grows as
$\#E^{1/6}$, literature attacks that are exhaustive search.

By family the record reads as follows. On \emph{Koblitz curves}
(\S\ref{sec:koblitz}) every September result below rho reversed against a
stronger rho, the $n=83$ online wins were retired by a matched
preprocessing rho, Table~A sits at $4.67$--$18.07\times$ (rule v8) or
$9.9$--$43\times$ (rule v4 recipe, canonical step), and every constructed or priced
route to ECC2K-130 is at least about 5 bits above rho (a hypothetical
genus-130 Jacobian would not be, but none is known). On \emph{random and
prime-degree binary curves} (\S\ref{sec:binary}) the solving degree of
Weil-descended systems grows with $\ell$ and tracks the semi-regular
prediction on every toy cell measured, while the first and last fall
degrees stay at 3--4; no algebraic oracle beats the pair table; and the five
validated Table~B rows are $10.6$--$38.2\times$ rho. On \emph{prime fields}
(\S\ref{sec:prime}) the eight validated Table~B rows are $5.5$--$13.5\times$
rho and rising, the Kummer tower's solving degree reaches 7 and is at least
8 at $N=21$, 24 rounds of structured bases on the P-256 field found no
relation for an unplanted target, and walks over a million curves isogenous
to P-256 fired no structural screen. The single-digit gaps all belong to
weak extension-field classes (\S\ref{sec:transfer}).

The honest summary is the campaign of \S\ref{sec:campaign}. On the
question the survey asks---does index calculus solve one unseen
elliptic-curve target, cold, for less than rho?---the answer on every curve
we measured, at every weight of the pricing band, is no. The closest arm is
a compact Frobenius-orbit Koblitz index calculus at $4$--$18\times$
signed-Frobenius rho in the campaign (separately, the leaderboard's Koblitz
pipeline is $9.9$--$43\times$ at 36.6--47.2 bits once rho is priced at the
canonical step), prime-field index calculus falls behind as $r^{1/6}$,
and no fitted exponent of an index-calculus arm lies below rho's. The
online $13.9\times$ at $n=61$ compared a table against no table and loses
$53$--$317\times$ to preprocessing rho with equal advice; it is retired.

Measured against its own baseline, which is the criterion the programme now
works to, index calculus keeps improving: about twenty changes landed on
14--15 October, among them a genus-3 index calculus after cubic GHS
whose exponent (about $q^{1.4}$) lies below rho's $q^{3/2}$, an $18\times$ cheaper
$k=4$ ladder and a further $18\times$ cheaper $k=4$ solve, and production
free relations on Koblitz curves ($0.960\times$). The genus-3 family is
the nearest to rho: $1.15\times$ a measured negation-map rho at $q=2^{20}$
after the second round; after rounds three and four, a preregistered
comparison put it below the strongest rho at $q=2^{16}$ at every band weight
under the registered pricing, unresolved at the other sizes to $2^{20}$,
and $2.7$--$12.4\times$ rho under a heavier sensitivity pricing of one
operation. That is the programme's only decision of index calculus below
the strongest rho; it is in a weak structured class, depends on that
pricing, and is not yet reviewed. The second is $k=4$ over
$\FF_{(2^m)^4}$, whose gap closes from $5{,}080\times$ to $6.1\times$ rho
between $n\approx2^{24}$ and $2^{56}$ at the $n^{-1/4}$ the theory
predicts; its extrapolated crossover near $2^{68}$ is the first in the
programme at a size where rho itself can be run, so it is the next
experiment. Both families are known to be weak in theory (Gaudry, Diem,
GHS); what the programme adds is a measured, verified approach to the
crossover, and neither is a deployed curve shape. The validation sweep of
19--20 October also retired the remaining online wins at $n=83$ and
reopened a replay regression in the root-finder session, since traced to an
x86-only code path in inherited F4.

The week to 24 October added a third family of the same kind: a
Joux--Vitse cover of $E(\FF_{p^6})$ (plain Nagao arm) whose gap closes as about $p^{-2}$ and
is $1.6$--$2.0\times$ a pooled rho constant at $p=1009$, with the crossover
extrapolated near $p\approx1330$; measuring rho and the cover's unpriced
work at that size is, with the $k=4$ run near $2^{68}$, a direct test of a
predicted crossover.
It also added four factor bases closed under known-logarithm maps, worth
$1.3$--$20\times$ against the programme's own arms and nothing in the
exponent, and a re-priced rho that nearly doubled the Koblitz leaderboard's
distance. The stages that now dominate are relation yield, factor-base size
and linear algebra. One thing the paper cannot yet say is how any of this
compares with published index calculus at the same parameters: every ratio
here is against rho or against the programme itself. Building that
reference arm (open problem~10) comes before any claim of improving on the
state of the art, and the other open problems of \S\ref{sec:open} say what
would move the stages that dominate.
